\section{Introduction}
Semantic segmentation is a fundamental task in computer vision, enabling pixel-level understanding of complex visual scenes. It has critical applications in autonomous driving, robotic manipulation, and aerial imagery analysis. Standard architectures, such as fully convolutional networks and vision transformers, have achieved remarkable accuracy. However, high-resolution images present a significant challenge. The computational cost of dense self-attention scales quadratically with the spatial dimensions of the feature maps, making them prohibitively expensive for real-time deployment on edge devices.

To mitigate this complexity, existing approaches often resort to downsampling input images or employing local attention mechanisms. While downsampling reduces memory overhead, it inevitably discards fine-grained spatial details, leading to blurred boundaries around small objects. Local attention, on the other hand, restricts the receptive field, preventing the model from capturing global contextual dependencies. This is particularly problematic for classes that require global spatial context to resolve ambiguity, such as large water bodies or continuous road surfaces in remote sensing \cite{miller2023aerial}.

In this paper, we propose SynapSeg, a contrastive multi-scale affinity learning framework that resolves these limitations. Our key insight is that pixel-level affinity can be routed sparsely across multiple scales, preserving long-range context without dense computation. We introduce a sparse affinity propagation module that learns to propagate context from low-resolution global features to high-resolution local features. Furthermore, we develop a localized contrastive loss to align these representations across resolutions, ensuring scale-invariant feature extraction.

We summarize our main contributions as follows:
\begin{enumerate}[label={\lbrack\roman*\rbrack}, leftmargin=*]
  \item We propose a sparse affinity propagation module that scales linearly with the image size, enabling long-range contextual integration \cite{wong2022sparse}.
  \item We introduce a localized contrastive loss that aligns affinity distributions across resolutions, mitigating scale-induced feature drift.
  \item We demonstrate state-of-the-art performance on multiple high-resolution benchmarks while reducing training memory footprint by $35\%$.
\end{enumerate}
